\documentclass[12pt, a4paper]{article}

% ── Packages ──────────────────────────────────────────────────────────────────
\usepackage[utf8]{inputenc}
\usepackage[T1]{fontenc}
\usepackage[margin=2.5cm]{geometry}
\usepackage{microtype}
\usepackage{parskip}
\usepackage{setspace}
\usepackage{titlesec}
\usepackage{titling}
\usepackage{fancyhdr}
\usepackage{booktabs}
\usepackage{array}
\usepackage{tabularx}
\usepackage[table]{xcolor}
\usepackage{hyperref}
\usepackage{enumitem}
\usepackage{lmodern}
\usepackage{csquotes}
\usepackage{mdframed}
\usepackage{ragged2e}

% ── Colors ────────────────────────────────────────────────────────────────────
\definecolor{darkblue}{HTML}{1a1a2e}
\definecolor{gold}{HTML}{c8a96e}
\definecolor{rust}{HTML}{8b4513}
\definecolor{lightgray}{HTML}{f5f0e8}
\definecolor{midgray}{HTML}{555555}
\definecolor{codebg}{HTML}{f0ede8}

% ── Hyperref setup ────────────────────────────────────────────────────────────
\hypersetup{
    colorlinks=true,
    linkcolor=rust,
    urlcolor=rust,
    citecolor=rust,
    pdftitle={KNOCK: Design Your Door — Final Report},
    pdfauthor={Mustafa Mert Can, Ali Fuat Akyemiş}
}

% ── Section title formatting ──────────────────────────────────────────────────
\titleformat{\section}
    {\large\bfseries\color{darkblue}}
    {\thesection.}{0.5em}{}
    [\vspace{2pt}\textcolor{gold}{\titlerule[0.5pt]}]

\titleformat{\subsection}
    {\normalsize\bfseries\color{rust}}
    {\thesubsection}{0.5em}{}

\titleformat{\subsubsection}
    {\normalsize\bfseries\color{darkblue}}
    {\thesubsubsection}{0.5em}{}

\titlespacing{\section}{0pt}{18pt}{6pt}
\titlespacing{\subsection}{0pt}{12pt}{4pt}
\titlespacing{\subsubsection}{0pt}{8pt}{3pt}

% ── Quote box ─────────────────────────────────────────────────────────────────
\newmdenv[
    leftmargin=2cm,
    rightmargin=2cm,
    innerleftmargin=12pt,
    innerrightmargin=12pt,
    innertopmargin=8pt,
    innerbottommargin=8pt,
    linecolor=gold,
    linewidth=0.5pt,
    backgroundcolor=lightgray,
    skipabove=8pt,
    skipbelow=8pt
]{quotebox}

% ── Code style ────────────────────────────────────────────────────────────────
\newmdenv[
    leftmargin=1cm,
    innerleftmargin=10pt,
    innerrightmargin=10pt,
    innertopmargin=6pt,
    innerbottommargin=6pt,
    linecolor=gold,
    linewidth=0.4pt,
    backgroundcolor=codebg,
    skipabove=6pt,
    skipbelow=6pt
]{codebox}

% ── Header / Footer ──────────────────────────────────────────────────────────
\pagestyle{fancy}
\fancyhf{}
\fancyhead[L]{\small\textcolor{midgray}{KNOCK: Design Your Door}}
\fancyhead[R]{\small\textcolor{midgray}{CSE 358 — Spring 2025–2026}}
\fancyfoot[C]{\small\textcolor{midgray}{\thepage}}
\renewcommand{\headrulewidth}{0.4pt}
\renewcommand{\headrule}{\hbox to\headwidth{\color{gold}\leaders\hrule height \headrulewidth\hfill}}
\setlength{\headheight}{25pt}
\addtolength{\topmargin}{-13pt}

% ─────────────────────────────────────────────────────────────────────────────
\begin{document}

% ── COVER PAGE ────────────────────────────────────────────────────────────────
\begin{titlepage}
    \centering
    \vspace*{3cm}

    {\Huge\bfseries\color{darkblue} KNOCK}\\[0.4em]
    {\large\itshape\color{rust} Design Your Door}\\[0.8em]
    \textcolor{gold}{\rule{\linewidth}{0.5pt}}\\[0.8em]
    {\small\color{midgray} CSE 358 Introduction to Artificial Intelligence}\\[0.2em]
    {\small\color{midgray} Final Project Report}\\[3cm]

    {\large\bfseries\color{darkblue} Mustafa Mert Can \quad|\quad Ali Fuat Akyemiş}\\[0.4em]
    {\small\color{midgray} 20220808705 \quad|\quad 20210808043}\\[0.6em]
    {\small\color{midgray} May 2026}\\[0.6em]
    {\small\ttfamily\color{rust}
        \href{https://github.com/AliFuatAkyemis/knock-design-your-door}%
             {github.com/AliFuatAkyemis/knock-design-your-door}}

    \vfill

    \begin{quotebox}
        \centering
        \itshape\color{rust}
        ``Mama, take this badge off of me / I can't use it anymore''\\[0.3em]
        \normalfont\small\color{midgray}
        --- Bob Dylan, \textit{Knockin' on Heaven's Door} (1973)
    \end{quotebox}
\end{titlepage}

\thispagestyle{empty}
\newpage

% ── ABSTRACT ─────────────────────────────────────────────────────────────────
\section*{Abstract}
\addcontentsline{toc}{section}{Abstract}
\textcolor{gold}{\rule{\linewidth}{0.5pt}}
\vspace{4pt}

\textbf{KNOCK: Design Your Door} is an AI-augmented hermeneutic system that analyzes
contemporary covers of Bob Dylan's 1973 meditation on mortality,
\textit{Knockin' on Heaven's Door.} By orchestrating three distinct generative AI
techniques --- Speech-to-Text (Groq Whisper), Large Language Model analysis
(Llama 3.3 70B), and Text-to-Image generation (Flux via Pollinations.ai) --- the
project constructs a multi-layered reinterpretation of any performer's relationship
to the song's themes. The result is a visual artifact generated in real time,
steeped in the aesthetic language of 1973: Kodachrome film grain, warm earth tones,
and imagery drawn from the Vietnam War counterculture. This report details the
technical architecture, artistic vision, and philosophical engagement behind the
system, culminating in a personal manifesto on mortality, interpretation, and the
act of knocking.

\newpage

% ── TABLE OF CONTENTS ─────────────────────────────────────────────────────────
\tableofcontents
\newpage

% ─────────────────────────────────────────────────────────────────────────────
\section{Introduction}

Bob Dylan wrote \textit{Knockin' on Heaven's Door} in 1973 for Sam Peckinpah's
revisionist Western \textit{Pat Garrett \& Billy the Kid.} The song accompanies
the death of a sheriff --- a man disarmed, exhausted, and surrendering to the
inevitable. In fewer than two minutes, Dylan distilled an entire era's grief into
four lines and a repeated refrain. The year was 1973: the Paris Peace Accords had
been signed in January, ostensibly ending direct American involvement in Vietnam,
yet the war's psychological wreckage endured. A generation raised on protest had
grown tired. The counterculture was fragmenting. To knock on heaven's door was not
merely to die --- it was to relinquish the burden of a world that had demanded too much.

Fifty years later, the song persists --- performed by hundreds of artists across
genres, languages, and continents. Each cover is a new interpretation: a contemporary
voice pressing its ear against Dylan's original door and hearing something different
behind it. KNOCK asks: what happens when we run that interpretation through an AI
pipeline designed to hear the song the way 1973 heard it?

The project is not a music generator, nor a simple chatbot. It is a
\textbf{hermeneutic machine} --- a system that interprets interpretations, and in
doing so, asks whether meaning can survive translation across time, medium,
and model weights.

% ─────────────────────────────────────────────────────────────────────────────
\section{Technical Implementation}

\subsection{System Architecture Overview}

KNOCK is implemented as a full-stack web application with a clear separation between
the AI processing backend and the thematic frontend. The pipeline is strictly
sequential: audio enters, passes through three AI stages, and a visual artifact exits.
The architecture is designed to be asynchronous throughout, ensuring the interface
remains responsive during heavy model inference.

\vspace{8pt}
\begin{center}
\renewcommand{\arraystretch}{1.4}
\begin{tabularx}{\linewidth}{
    >{\bfseries\color{darkblue}}l
    >{\ttfamily\small}l
    >{\RaggedRight\arraybackslash}X
}
\rowcolor{darkblue}
\textcolor{white}{\textbf{Layer}} &
\textcolor{white}{\textbf{Technology}} &
\textcolor{white}{\textbf{Role}} \\
\rowcolor{lightgray}
Frontend         & Vanilla JS + CSS                  & Audio recording/upload, thematic UI, film-grain overlays \\
Backend API      & FastAPI (Python)                  & Async routing, service orchestration, error handling \\
\rowcolor{lightgray}
Stage 1 --- STT  & Groq Whisper (large-v3-turbo)     & Audio transcription \\
Stage 2 --- LLM  & Groq Llama 3.3 70B                & Philosophical and thematic analysis \\
\rowcolor{lightgray}
Stage 3 --- T2I  & Pollinations.ai Flux              & Visual artifact generation \\
Prompt Layer     & utils/prompt\_engineer.py         & 5-layer 1973 aesthetic injection \\
\end{tabularx}
\end{center}
\vspace{8pt}

\subsection{AI Technique Integration}

The three AI techniques are not simply chained in sequence --- each stage's output
actively shapes the next stage's input, creating emergent results that no single
technique could produce alone.

\subsubsection{Stage 1 --- Speech-to-Text (Groq Whisper)}

The \textit{whisper-large-v3-turbo} model transcribes the user's audio cover with
high fidelity, capturing not just the words but the performer's unique phrasing,
emphasis, and vocal dynamics encoded in the text representation. This transcription
is passed verbatim to Stage 2, ensuring that the LLM analysis is grounded in the
specific human performance rather than a generic lyric sheet.

\subsubsection{Stage 2 --- Large Language Model Analysis (Llama 3.3 70B)}

The LLM receives the Whisper transcription along with a carefully engineered system
prompt that grounds it in the 1973 historical context. The model is instructed to
analyze: (1) the emotional register of the performance, (2) connections to the
Vietnam War era and counterculture themes, (3) the performer's relationship to the
song's core motifs --- farewell, mortality, disarmament --- and (4) a symbolic visual
direction for Stage 3. This analysis is structured JSON output, which the backend
parses and routes to the prompt engineering layer.

\subsubsection{Stage 3 --- Text-to-Image Generation (Pollinations.ai Flux)}

Flux receives a dynamically constructed prompt from \texttt{utils/prompt\_engineer.py},
which synthesizes the LLM's thematic analysis with the 5-layer aesthetic framework
described in Section~\ref{sec:prompt}. The resulting image is the project's primary
artistic output: a visual artifact that is both historically anchored and uniquely
responsive to the specific human performance that triggered it.

\subsection{Technical Ambition: The 5-Layer Prompt Engineering System}
\label{sec:prompt}

The prompt engineering layer is the technical heart of KNOCK's artistic ambition.
Rather than passing the LLM's analysis directly to the image model,
\texttt{utils/prompt\_engineer.py} processes it through five injection layers:

\begin{itemize}[leftmargin=1.5em, itemsep=4pt]
    \item \textbf{Aesthetic Base:} Every prompt includes Kodachrome film grain,
    35mm documentary photography style, shallow depth of field, warm earth tones,
    and the faded yellows and burnt oranges characteristic of 1973 photojournalism.

    \item \textbf{Historical Layer:} Keywords from the LLM analysis trigger
    period-specific imagery: dog tags, protest banners, military helicopters,
    rice paddies, candlelight vigils, or the smoky backrooms of American folk clubs.

    \item \textbf{Symbolic Anchor:} A `Door' motif --- weathered wood, peeling paint,
    threshold light --- is injected as a compositional constant, providing visual
    continuity across all generated artifacts.

    \item \textbf{Emotional Modulator:} The LLM's extracted emotional register
    (grief, resignation, defiance, hope) adjusts the prompt's descriptive language:
    a grief-coded performance generates heavier shadow and empty space; a defiant one
    generates wider aperture and stronger contrast.

    \item \textbf{Negative Prompt Layer:} Explicit exclusions prevent the model from
    generating anachronistic imagery, digital aesthetics, or saturated colors
    incompatible with the 1973 visual language.
\end{itemize}

\subsection{Repository Structure}

The full source code and project assets are organized as follows:

\begin{codebox}
\small\ttfamily
KNOCK-PROJECT/\\
|-- backend/                \textcolor{midgray}{\# FastAPI server and AI logic}\\
|   |-- config.py           \textcolor{midgray}{\# Environment configuration}\\
|   |-- main.py             \textcolor{midgray}{\# Application entry point}\\
|   |-- models/             \textcolor{midgray}{\# Pydantic data schemas}\\
|   |   `-- schemas.py\\
|   |-- services/           \textcolor{midgray}{\# AI Model integrations}\\
|   |   |-- gemini\_service.py\\
|   |   |-- siliconflow\_service.py\\
|   |   `-- whisper\_service.py\\
|   `-- utils/              \textcolor{midgray}{\# Helper utilities}\\
|       `-- prompt\_engineer.py \textcolor{midgray}{\# 5-layer prompt system}\\
|-- frontend/               \textcolor{midgray}{\# Client-side implementation}\\
|   |-- app.js              \textcolor{midgray}{\# Audio handling and API calls}\\
|   |-- index.html          \textcolor{midgray}{\# Thematic user interface}\\
|   `-- style.css           \textcolor{midgray}{\# 1973 film-grain aesthetics}\\
|-- .env                    \textcolor{midgray}{\# Secret API keys}\\
`-- requirements.txt        \textcolor{midgray}{\# Python dependencies}
\end{codebox}

\subsection{Setup and Execution Walkthrough}

Since sensitive configuration files (\texttt{.env}) and local environments (\texttt{venv/}) are excluded from version control for security and efficiency, follow these steps to initialize and run the system after cloning:

\subsubsection{1. Clone the Repository}
Clone the project from GitHub to your local machine:
\begin{codebox}
\small\ttfamily
git clone https://github.com/AliFuatAkyemis/knock-design-your-door.git\\
cd knock-design-your-door
\end{codebox}

\subsubsection{2. Create Environment Configuration}
The \texttt{.env} file is not included in the repository to protect API secrets. You must create it manually in the root directory:
\begin{codebox}
\small\ttfamily
touch .env \textcolor{midgray}{\# Or create a new file named .env using a text editor}
\end{codebox}
Open the \texttt{.env} file and insert the configuration. Note that while most variables have defaults, the \texttt{GROQ\_API\_KEY} is strictly required:
\begin{codebox}
\small\ttfamily
GROQ\_API\_KEY=your\_groq\_api\_key\_here\\
MAX\_UPLOAD\_SIZE\_MB=25\\
UPLOAD\_DIR=uploads\\
OUTPUT\_DIR=outputs\\
FRONTEND\_ORIGIN=http://localhost:8000
\end{codebox}

\subsubsection{3. Initialize Virtual Environment \& Install Dependencies}
Create a fresh virtual environment to manage project dependencies:
\begin{codebox}
\small\ttfamily
python -m venv venv\\
source venv/bin/activate  \textcolor{midgray}{\# On Windows: venv\textbackslash Scripts\textbackslash activate}\\
pip install -r requirements.txt
\end{codebox}

\subsubsection{4. Launch the Server}
Start the FastAPI backend using Uvicorn:
\begin{codebox}
\small\ttfamily
uvicorn backend.main:app --reload --port 8000
\end{codebox}

Once the server is running, navigate to \texttt{http://localhost:8000} in your web browser to access the KNOCK interface.

% ─────────────────────────────────────────────────────────────────────────────
\section{Artistic Originality}

\subsection{The Hermeneutic Chain}

The conceptual backbone of KNOCK is what we call the \textbf{Hermeneutic Chain}:
a sequence of interpretations in which each link transforms meaning.

\begin{quotebox}
    \centering\itshape\color{rust}
    Dylan (1973) $\;\to\;$ Cover Artist $\;\to\;$ Groq Whisper $\;\to\;$
    Llama 3.3 $\;\to\;$ Flux $\;\to\;$ Viewer
\end{quotebox}

Dylan encoded grief, farewell, and the exhaustion of idealism into a two-minute folk
song. A contemporary performer encodes their own mortality, their own relationship to
loss, into that song's shell. Whisper encodes the acoustic performance into text.
Llama encodes the text into thematic meaning. Flux encodes that meaning into pixels.
And the viewer encodes all of it into personal experience.

This is not degradation --- it is amplification. Each translation adds a layer of
human and machine subjectivity, and the final artifact carries the sediment of all
of them. The viewer does not see what Dylan wrote; they see what it means, filtered
through five interpreters, to look at a door and not yet know what lies behind it.

\subsection{Aesthetic Coherence}

Every design decision in KNOCK --- from the frontend's desaturated color palette to
the film-grain CSS filter applied to generated images --- serves the single aesthetic
proposition: \textit{1973 is not a date, it is a feeling.} The feeling of a world
that had fought for something, had almost believed in it, and was now slowly setting
it down.

The frontend deliberately avoids modern UI conventions: no gradients, no animations
beyond a slow dissolve, no sans-serif clarity. The interface feels like a document
found in a drawer. The generated images feel like photographs from a war that ended
before anyone agreed on what it meant.

\subsection{Emotional Architecture}

KNOCK is designed to produce a specific emotional sequence in the user:
\textbf{participation $\to$ surrender $\to$ reflection.} The user participates by
recording or uploading their cover --- an intimate, slightly vulnerable act. They
surrender control to the pipeline, which interprets them without negotiation. They
reflect when the artifact appears: an image that is about them, generated by a
machine, rooted in a world before they were born.

This sequence mirrors the song itself: the sheriff removes his badge (participation),
accepts death (surrender), and becomes part of something larger than himself
(reflection). The user does not die; but they do, briefly, let the machine be the
one doing the listening.

% ─────────────────────────────────────────────────────────────────────────────
\section{The Artwork}

\subsection{Visual Language}

The generated artifact --- what we call the \textbf{Artifact} --- is not an
illustration of the song. It is a diagnostic portrait of the performer's relationship
to the song. Two people performing identical lyrics will generate different artifacts,
because Whisper captures how they sing, the LLM reads what that reveals about their
inner state, and Flux renders that inner state in 1973's visual language.

The Kodachrome palette was chosen with deliberateness: it is the palette of Eddie
Adams' 1968 photograph of the Saigon execution, of Nick Ut's 1972 `Napalm Girl,'
of the documentary footage that brought the Vietnam War into American living rooms.
These images did not need to be beautiful to be powerful; they needed to be present.
KNOCK's artifacts inherit that obligation.

\subsection{The Door Motif}

The Door appears in every artifact as a compositional anchor --- sometimes as a
literal weathered door in an adobe wall, sometimes as a threshold of light at the
edge of darkness, sometimes as the silhouette of a figure standing at a gate. It is
the one element the prompt engineering layer protects from variation. Everything else
--- the historical details, the emotional weather, the figure in the foreground ---
responds to the specific performer. The Door does not move. The Door waits.

% ─────────────────────────────────────────────────────────────────────────────
\section{Philosophical Engagement}

\subsection{Historical Depth: 1973 and the Grammar of Farewell}

To understand \textit{Knockin' on Heaven's Door,} one must understand what 1973
meant to carry. The Paris Peace Accords were signed on January 27, 1973, ending
direct American military involvement in Vietnam --- but not ending the war, not
healing the 58,000 Americans and roughly two million Vietnamese dead, and not
resolving the fracture the conflict had opened in American consciousness. The
Watergate hearings began in May. A country that had organized itself around
opposition to a war suddenly had to figure out what it was for.

Dylan, who had been the voice of the protest generation in the early 1960s, had
spent the intervening years deliberately retreating from that role. By 1973, he was
in Durango, Mexico, playing a minor role in Peckinpah's elegiac Western and writing
a song for a dying man. The sheriff in the film is a figure from an older moral
order --- a man whose codes no longer apply. His death is not heroic; it is a
laying-down. Dylan, who had once written ``The times they are a-changin','' now
wrote about the relief of no longer having to change anything.

KNOCK's prompt engineering layer was designed to carry this specific historical
weight. The LLM system prompt includes a carefully researched paragraph on the 1973
context, ensuring that even a contemporary pop cover of the song is analyzed through
the lens of this specific farewell, rather than the generic sentimentality the song
has accumulated over five decades of use.

\subsection{Dylan's Worldview: The Ethics of Withdrawal}

Dylan's 2016 Nobel Prize lecture offers a retrospective key to his 1973 sensibility.
Speaking of \textit{Moby Dick,} he described the captain's obsession as a kind of
moral insanity --- the refusal to let go. The sheriff in \textit{Pat Garrett \&
Billy the Kid} is the opposite: a man who has the wisdom to put down his badge,
to acknowledge that his authority has become more burden than purpose. Dylan's Nobel
lecture suggests he has always been interested in this figure --- the one who stops,
who refuses to keep fighting a war that has already been decided.

This is the worldview KNOCK inherits. The system does not celebrate death; it
explores the specific dignity of laying something down. When the LLM analyzes a
performer's vocal delivery, it is trained to look for this: are they carrying
something? And can you hear them, in this song, beginning to set it down?

\subsection{The AI as Interpretive Subject}

A central philosophical question animates this project: can an algorithmic system
genuinely \textit{engage} with a work of art, or does it merely process its surface?
We do not claim that Llama 3.3 feels grief. We do claim something more modest and
more interesting: that the model operates within a latent space of human cultural
production, and that its outputs, when properly constrained, can reflect the
structure of human meaning-making rather than its mere content.

The LLM has been trained on millions of texts about mortality, war, folk music, and
loss. It cannot feel these things, but it has learned the grammar of how humans speak
about them. When it analyzes a performer's transcription and identifies ``a quality
of resignation threaded through the second verse, consistent with the 1973
counterculture's exhausted acceptance of systemic failure,'' it is doing something
real --- not feeling, but pattern-matching against the deepest patterns of human
expression. We find this sufficient, and more honest than claiming less.

% ─────────────────────────────────────────────────────────────────────────────
\newpage
\section{Artist's Manifesto}

\begin{quotebox}
    \centering\itshape\color{rust}
    ``Everyone knocks on heaven's door in their own way. What matters is not what
    you find behind it, but who you discover yourself to be while knocking.''
\end{quotebox}

\subsection{Why This Medium?}

We chose this medium --- a live processing pipeline that generates images from human
singing --- because it is the only form we could imagine that would force the question
we wanted to ask. The question is not ``what does this song mean?'' That question
has been answered ten thousand times, in ten thousand key changes, by everyone from
Eric Clapton to Guns N' Roses to Arlo Guthrie's daughter. The question we wanted to
ask is harder: \textit{what does it mean to interpret something that has already been
interpreted until it barely knows its own name?}

A generative image cannot answer that question in isolation. An LLM cannot answer it
alone. But a pipeline that takes a human voice --- wavering, specific, carrying its
own history --- and runs it through layers of machine interpretation before producing
a visual artifact: that pipeline makes the question visible. You can see the
interpretation happening. You can watch the song travel.

We chose this medium also because it requires the user to do something slightly
uncomfortable: to sing. Not to type about the song, not to click through it, but to
open their mouth and make the sounds. This is the oldest form of engagement with folk
music, and we wanted to preserve it even inside a machine. The Whisper model does not
care if your voice is good. The LLM does not judge your pitch. But you have to
decide, in the moment of recording, how to hold these particular words. That decision
is already an interpretation. That is where the art begins --- before the AI has done
anything at all.

The web application format was chosen because it removes barriers. A gallery
installation would have been more spatially dramatic, but it would have restricted
access. A command-line tool would have been more technically pure, but it would have
excluded everyone who doesn't live in a terminal. The browser is where most people
encounter things now. We wanted KNOCK to be available to anyone who has ever hummed
this song in a car and wondered what they were really singing about.

\subsection{What Caught Us?}

We encountered this song the way most people our age do: late, mediated, already
filtered through decades of use. We did not hear it in 1973, in a dark cinema in the
aftermath of Vietnam. We heard it covered, sampled, placed in television soundtracks,
used as a backdrop for sports montages and political advertisements. The song had been
polished smooth by repetition. We thought we knew what it was.

What caught us --- genuinely, in a way that changed how we thought about the project
--- was watching the Peckinpah film. Specifically, the scene for which Dylan wrote
the song: Slim Pickens, playing a dying sheriff, lying in a field by a river while
his wife holds his hand. It is not a dramatic death. There is no speech, no
revelation, no final stand. He simply lies there, and the river moves, and the light
changes, and Dylan sings those four lines. And something in that image --- a man at a
threshold, not crossing it dramatically but simply arriving at it, spent --- felt
true in a way we had not expected.

We had come to the project thinking about AI and music and generative systems. We
left it thinking about what we carry, and what we might one day want to set down.
That shift is, we think, what the assignment was asking for. Not a technical
demonstration, but an encounter. We had the encounter. It arrived through a
sixty-year-old actor dying in a Sam Peckinpah film, and we did not see it coming.

We were also caught by the song's refusal of specificity. Dylan does not tell us what
heaven is, or whether it exists, or whether the door actually opens. He tells us the
badge is heavy, the guns are too dark to be clean, and there is a knock. That is all.
The song's power is entirely in what it withholds. This restraint became a design
principle for KNOCK: we show the door, we do not show what is behind it. The artifact
is always a threshold image, never a revelation. We earned that choice from Dylan.

\subsection{AI as What?}

This is the question we debated longest. The easy answers came quickly and were
quickly discarded.

\textit{AI as tool?} Too simple. A hammer is a tool; it has no relationship to what
you build with it. Llama 3.3 has read more about mortality and folk music and the
Vietnam War than either of us will in our lifetimes. When it analyzes a performance,
it brings that accumulated reading to bear. It is not neutral. Calling it a tool
flattens something real.

\textit{AI as collaborator?} Too generous, or too specific. A collaborator has
intentions, preferences, the capacity to surprise you on purpose. The LLM surprised
us constantly --- producing analyses we would not have written, making connections
between a performer's phrasing and specific historical events that we had not
anticipated. But it did not do this \textit{intentionally.} It did it because it was
well-trained and the prompt was well-constructed. The surprise was real; the intention
was ours.

\textit{AI as mirror?} Closer. When KNOCK analyzes a cover performance, it reflects
back a reading of that performance that the performer did not consciously generate.
It reveals the subtext. A user who sings the song with forced brightness will generate
an artifact with shadows the user did not choose. The system sees what the user
performs, not what the user intends. That is mirror-like behavior.

Our final answer: \textbf{AI as archaeologist.} The models we use --- Whisper, Llama,
Flux --- were trained on enormous archives of human production: speech, text, images.
They are, in a meaningful sense, compressed human culture. When they process a
performance of a 1973 folk song, they are not computing from scratch; they are
excavating from within a vast sediment of human expression, finding the layer that
corresponds to grief, to farewell, to the specific visual grammar of that era. They
dig, and what they bring up is genuinely from the earth --- the human earth --- even
if the hands doing the digging are algorithmic.

This metaphor shaped our technical decisions. We did not try to make the AI create
something `new' in the sense of unprecedented. We tried to make it excavate something
specific: the 1973 layer. Every element of the prompt engineering system is a
constraint that forces the excavation to go in the right direction. The art is in
knowing where to dig.

\subsection{Your Door: On Thresholds We Carry}

This question --- what does knocking on heaven's door mean in your own life? --- is
the one we approached most carefully, because it is the one where honesty is most
expensive.

We are students of computer science in a moment of profound transition. The field we
are entering is not the field it was five years ago, and it will not be the same field
in five years more. The tools we are learning to build are tools that are, in some
configurations, learning to replace the need for builders. We are training for a
profession that is itself being retrained. This is our version of 1973: the feeling
that the old frameworks are dissolving, that the badge --- the credential, the
discipline, the role --- may soon be too heavy to wear, or may mean something
different than it used to.

We do not say this with self-pity. The sheriff in the film does not die with
self-pity. He dies with the specific dignity of a man who has chosen to stop fighting
something larger than himself and instead attend to the river, the light, his wife's
hand. There is something instructive in that choice. We do not yet know what it means
to stop fighting the transitions of our field and simply attend to what is good in
the work --- but KNOCK is, among other things, an attempt to practice that attention.

We also encountered personal thresholds in building this project. Working together
across late nights and debugging sessions, we discovered that collaboration has its
own doors --- moments where you must decide whether to hold your position or let your
partner's vision lead. Those moments felt, in small ways, like the song: the relief
of releasing something you had been gripping too hard.

What does knocking on heaven's door mean to us? It means: there are things you cannot
force. There are doors that open on their own schedule, or not at all. The knocker's
job is only to knock --- to keep showing up at the threshold, with whatever you have,
and wait. KNOCK, the system, does this literally: it takes a human at a threshold
(singing a song about thresholds) and renders that threshold back to them in visual
form. We designed it this way because we believe the rendering is the point.
Not the arrival. Never the arrival. Only the knocking.

We will graduate from this program and walk through a door that none of us can see
clearly yet. The field is changing; the economy is changing; what it means to build
software, to make something, to be useful --- all of it is in motion. We made this
project as an act of attention to that threshold. Not an answer. An attention.
We hope it is the kind of attention that leaves a mark.

\vspace{12pt}
\begin{mdframed}[linecolor=gold, linewidth=0.4pt, backgroundcolor=lightgray,
    innerleftmargin=10pt, innerrightmargin=10pt,
    innertopmargin=8pt, innerbottommargin=8pt]
\small\itshape\color{midgray}
\textbf{Team contribution note:} Technical architecture, backend implementation, and
prompt engineering system were developed jointly with continuous iteration between
both team members. Artistic direction and manifesto writing were co-authored through
a deliberate process of challenge and revision. Both members are prepared to discuss
any aspect of the project in full during the exhibition.
\end{mdframed}

% ─────────────────────────────────────────────────────────────────────────────
\section{Conclusion}

KNOCK: Design Your Door succeeds in doing what we set out to do: building a system
that does not merely process a song, but engages with it. The technical pipeline ---
Whisper to Llama to Flux, mediated by a carefully engineered prompt layer ---
produces artifacts that are visually distinctive, historically grounded, and genuinely
responsive to individual human performance. No two artifacts are the same, because
no two performances are the same.

More importantly, the system taught us something we did not know we needed to learn.
We came to it as engineers interested in pipeline architecture. We left it as people
who had spent weeks thinking about farewell, about the dignity of disarmament, about
what it means to stand at a threshold and knock without knowing what --- or who ---
might answer.

Dylan wrote a song for a dying man in a film about a dying era. We built a machine
for listening to how people sing that song today, and for rendering what we hear in
the visual grammar of that era. The door in every artifact is the same door. The
light behind it changes every time. That variability is not a bug. It is the entire
point.

\begin{quotebox}
    \centering\itshape\color{rust}
    ``Knock. The door is yours to design. What's behind it is yours to discover.''
\end{quotebox}

% ─────────────────────────────────────────────────────────────────────────────
\section{Academic Integrity Statement}

All AI tools used in this project are documented transparently. The Groq API was used
for Whisper (\texttt{whisper-large-v3-turbo}) and Llama 3.3
(\texttt{llama-3.3-70b-versatile}) inference. The Pollinations.ai API was used for
Flux-based image generation. All creative vision, architectural design, prompt
engineering decisions, and philosophical reflection are original work by the project
team. No code was borrowed without attribution; the full repository is available for
inspection at
\href{https://github.com/AliFuatAkyemis/knock-design-your-door}%
     {\texttt{github.com/AliFuatAkyemis/knock-design-your-door}}.
Both team members can explain every component of the pipeline in full.

\end{document}
